% Бие даалт 2 постер (нэг нүүр, A3) - F.CSM302 Үйлдлийн систем, баг II-4-12 (Producer-Consumer).
% One font (Helvetica Neue), three colours (ink, produce blue, consume purple), real screenshots only.
%   cd docs/latex && tectonic -o build II-4-12_Poster.tex
\documentclass{article}
\usepackage[a3paper,margin=13mm]{geometry}
\usepackage{fontspec}
\setmainfont{Helvetica Neue}[UprightFont=*, BoldFont=* Bold, ItalicFont=* Italic, BoldItalicFont=* Bold Italic]
\setsansfont{Helvetica Neue}[UprightFont=*, BoldFont=* Bold, ItalicFont=* Italic, BoldItalicFont=* Bold Italic]
\setmonofont{Helvetica Neue}[UprightFont=*, BoldFont=* Bold]
\usepackage{graphicx}
\usepackage[table]{xcolor}
\usepackage{tikz}
\usetikzlibrary{arrows.meta,positioning}
\usepackage[most]{tcolorbox}
\usepackage{qrcode}
\usepackage{array}
\usepackage{enumitem}
\graphicspath{{../screenshots/}}
\pagestyle{empty}
\setlength{\parindent}{0pt}
\hyphenpenalty=10000 \exhyphenpenalty=10000
\linespread{1.08}

\definecolor{ink}{HTML}{111827}
\definecolor{produce}{HTML}{4F7CF7}
\definecolor{consume}{HTML}{8B5CF6}
\definecolor{soft}{HTML}{F3F5FB}
\definecolor{edge}{HTML}{D5DAE5}
\color{ink}

\newcommand{\Release}{https://github.com/khuslenshi/os-bd1-producer-consumer/releases/tag/v1.1.0}

% A block: coloured title strip + white body
\newtcolorbox{block}[2]{enhanced, colback=white, colframe=edge, boxrule=0.8pt, arc=6pt,
  left=10pt, right=10pt, top=6pt, bottom=6pt, fonttitle=\bfseries\fontsize{16}{19}\selectfont,
  coltitle=white, colbacktitle=#1, title=#2, toptitle=5pt, bottomtitle=5pt, before skip=0pt, after skip=6pt,
  before upper=\raggedright}
% A real screenshot in a thin frame with a caption
\newcommand{\shot}[4]{{\setlength{\fboxsep}{0pt}\setlength{\fboxrule}{0.8pt}%
  \centering{\color{edge}\fbox{\includegraphics[width=#1,trim=#3,clip]{#2}}}\par\vspace{4pt}
  {\fontsize{11}{13.5}\selectfont #4\par}}}
\newcommand{\metric}[3]{\begin{tcolorbox}[enhanced, colback=soft, colframe=soft, arc=6pt, boxrule=0pt,
  left=6pt, right=6pt, top=6pt, bottom=6pt, halign=center, before skip=0pt, after skip=0pt, width=\linewidth]
  {\color{#3}\bfseries\fontsize{26}{28}\selectfont #1}\par\vspace{2pt}{\fontsize{11}{13}\selectfont #2}\end{tcolorbox}}
\setlist{leftmargin=1.3em, itemsep=3pt, topsep=2pt, parsep=0pt}

\begin{document}
\fontsize{12}{15}\selectfont

% ================================================================ header
\begin{tcolorbox}[enhanced, colback=ink, colframe=ink, arc=8pt, left=16pt, right=16pt, top=9pt, bottom=9pt, after skip=7pt]
\color{white}
\begin{minipage}[c]{0.8\linewidth}\raggedright
{\fontsize{12.5}{15}\selectfont F.CSM302 Үйлдлийн систем · Бие даалт 2 · Баг \textbf{II-4-12}\par}
\vspace{5pt}
{\bfseries\fontsize{32}{36}\selectfont The Producer-Consumer Problem\par}
\vspace{5pt}
{\fontsize{15}{18}\selectfont Бодит C thread-үүдийн ажиллагааг "Buffer animation"-аар харуулах вэб апп\par}
\vspace{7pt}
{\fontsize{12}{14.5}\selectfont А.Эрхэмтөр B241940029 · Б.Хүслэн B242270107 · Н.Батмөрөөдөл B231910039 · Г.Оюунгэрэл B232270802 · С.Итгэл B242270031\par}
\end{minipage}\hfill
\begin{minipage}[c]{0.18\linewidth}\raggedleft
{\fontsize{12}{14.5}\selectfont Сурах бичиг: ch.7,\\ Programming Project 4\\[4pt] Хувилбар \textbf{v1.1.0}\par}
\end{minipage}
\end{tcolorbox}

% ================================================================ left column
\begin{minipage}[t]{0.488\linewidth}\vspace{0pt}% top-align the two columns
\begin{block}{produce}{Асуудал}
Хэд хэдэн \textbf{producer} нэг дундын хязгаартай буферт item тавьж, \textbf{consumer} авна.
\begin{itemize}
  \item Буфер дүүрэн бол producer хүлээх ёстой.
  \item Хоосон бол consumer хүлээх ёстой.
  \item Хоёр thread буферийг зэрэг өөрчилбөл өгөгдөл алдагдана: \textbf{race condition}.
\end{itemize}
Thread-ийн ажиллагаа нүдэнд харагддаггүй тул шийдлийг бодит программ дээр нь харуулна.
\end{block}

\begin{block}{produce}{Шийдэл}
Номын дарааллаар: \textbf{wait(empty)}, \textbf{wait(mutex)}, тавих, \textbf{signal(mutex)}, \textbf{signal(full)}. Consumer тэгш хэмтэй. Программ үйлдэл бүрийг JSON мөрөөр хэвлэж, дэлгэц зөвхөн харуулна.\par\vspace{7pt}
\centering
\begin{tikzpicture}[font=\fontsize{11}{13}\selectfont, node distance=5mm,
  box/.style={draw=edge, line width=0.8pt, rounded corners=5pt, fill=soft, align=center, inner sep=5pt, text width=3.25cm},
  arr/.style={-{Stealth[length=2.4mm]}, line width=1pt, color=produce}]
  \node[box] (core) {\textbf{C core}\\pthreads, semaphore\\\textbf{empty}, \textbf{full}, mutex};
  \node[box, right=of core] (api) {\textbf{/api/run}\\core-ийг ажиллуулж\\мөрийг stream хийнэ};
  \node[box, right=of api] (ui) {\textbf{Browser}\\буфер, thread карт,\\timeline, event log};
  \draw[arr] (core) -- (api);
  \draw[arr] (api) -- (ui);
\end{tikzpicture}
\end{block}

\begin{block}{produce}{Туршилт ба дүрслэл}
\textbf{Өөрчилж болно:} producer, consumer (0-16), буфер (1-16), унтах хугацаа, ажиллах хугацаа, seed, mutex асаах/унтраах.\par\vspace{5pt}
\textbf{Ажиглана:}
\begin{itemize}
  \item шар карт: semaphore хүлээж байна (empty = 0 эсвэл full = 0);
  \item хар хүрээ: mutex эзэмшигч, нэг удаад нэг;
  \item улаан: race, "count = X, Y байх ёстой";
  \item төгсгөлд Invariant ба Item шалгалт.
\end{itemize}
Timeline, event log, 0.1-1× хурд, алхамчлах, 12 алхамтай дуут хичээл, асуулттай 3 түвшин.
\end{block}

\begin{block}{produce}{Үр дүн}
\begin{minipage}[t]{0.485\linewidth}\metric{20/20}{core тест (Linux), 19/19 macOS}{produce}\end{minipage}\hfill
\begin{minipage}[t]{0.485\linewidth}\metric{0}{ThreadSanitizer анхааруулга, mutex-тэй (өмнө 51)}{produce}\end{minipage}\par\vspace{7pt}
\begin{minipage}[t]{0.485\linewidth}\metric{27/30}{mutex-гүй run-д race илэрсэн (өмнө 20/30)}{consume}\end{minipage}\hfill
\begin{minipage}[t]{0.485\linewidth}\metric{30/30}{mutex-тэй run-д алдагдсан 0, давхар 0}{consume}\end{minipage}\par\vspace{7pt}
{\fontsize{11}{13.5}\selectfont CI: Ubuntu (Node 22), macOS (Node 24), ThreadSanitizer job, \textbf{npm audit 0}. Гүйцэтгэл өөрчлөгдөөгүй (≈80 000 үйлдэл/с).\par}
\end{block}
\end{minipage}\hfill
% ================================================================ right column
\begin{minipage}[t]{0.488\linewidth}\vspace{0pt}% top-align the two columns
\begin{block}{consume}{Демонстрац: бодит программын дэлгэц}
\shot{0.9\linewidth}{web_02_buffer_full.png}{320 1200 320 500}{\textbf{Буфер дүүрэх} (4 producer, 1 consumer, буфер 2): \textbf{empty = 0} тул producer-ууд шар болж хүлээж байна.}
\vspace{6pt}
\shot{0.9\linewidth}{web_04_race.png}{320 1120 360 690}{\textbf{Race condition} (8 + 8 thread, mutex-гүй): P8 тавихад count = 2 байх ёстой атлаа 1.}
\vspace{6pt}
\shot{\linewidth}{web_12_items.png}{560 1080 560 500}{\textbf{Төгсгөлийн шалгалт}: item бүр дугаартай тул 185 item алдагдсан, 186 нь давхар авагдсаныг илрүүлэв.}
\end{block}

\begin{block}{consume}{Шүүмжийн дагуу хийсэн сайжруулалт}
II-4-7, V-1-8 багийн 14 санал: 10 зөвшөөрсөн, 3 хэсэгчлэн, 1 үгүй.
\begin{enumerate}[leftmargin=1.4em]
  \item \textbf{Race-ийг item-ээр илрүүлэх} (Major): тоолуур таарсан ч алдагдсан, давхар item илэрнэ.
  \item \textbf{Data race 0}: log-ийн уншилт, зогсоох флаг atomic, CI-д ThreadSanitizer.
  \item \textbf{Аюулгүй dependency}: Next.js 16.4.4, critical эмзэг байдал 0, Origin шалгалт.
  \item \textbf{Manual}: Linux дээр Node суулгах, UI ба core-ийн хүрээ, Node 22.18+.
\end{enumerate}
\end{block}

\begin{block}{produce}{Дүгнэлт ба GitHub}
Номын шийдлийг бодит thread-ээр ажиллуулж, алхам бүрийг харуулсан. Mutex асаалттай бүх run зөв, унтраавал race тогтмол илэрнэ. Хязгаар: Windows дээр WSL2 хэрэгтэй, race магадлалтай.\par\vspace{7pt}
\begin{minipage}[c]{0.34\linewidth}\qrcode[height=2.9cm]{\Release}\end{minipage}\hfill
\begin{minipage}[c]{0.62\linewidth}\raggedright
{\fontsize{11.5}{14}\selectfont\textbf{Final release v1.1.0}\\ github.com/khuslenshi/\\os-bd1-producer-consumer\\[4pt] \textbf{npm install}, \textbf{npm run dev}\\ README: Quick Start\par}
\end{minipage}
\end{block}
\end{minipage}

\end{document}
